\clearpage
\section{Appendix: Liturgical Resolution}
The celebration is the Twenty-ninth Sunday in Ordinary Time, Year A, on 18 October 2026 in the dioceses of the United States: green, Lectionary no.145, formula PC-S55-A. The controlling Missal is the \work{Roman Missal, Third Edition}, U.S. English (2011), based on the third typical Latin edition, emended reprint (2008). The controlling reading book is the U.S. \work{Lectionary for Mass}, second typical edition, volume I. The national calendar appoints this Sunday on the requested date; the occurring fixed feast of Saint Luke supplies no added proper here. No particular diocesan or religious-calendar overlay was specified.

The national calendar's Year A agrees with the requested study. Its weekday Cycle II is a separate cycle and changes none of these Sunday readings. Week XXIX owns the common Missal prayers and antiphons across the Sunday cycles; Year A owns the present reading selection. The Old Testament/Gospel correlation and the independently semicontinuous apostolic course follow the General Introduction to the Lectionary, §§65--69,103--104.

\begin{branchtable}
\bid{entrance-missal-or-permitted-chant} & GIRM 48; chant and recitation rules & permitted & Entrance & Missal route studied; no local chant selected.\\
\bid{sprinkling-in-place-of-penitential-act} & GIRM 51; Sunday option & permitted & Ordinary & Unselected; no sprinkling text supplied.\\
\bid{psalm-lectionary-or-permitted-sung-form} & GIRM 61 & permitted & Psalm & Lectionary route studied, with its actual strophes.\\
\bid{acclamation-lectionary-or-graduale} & GIRM 62--63 & permitted & Acclamation & Lectionary route studied; no other selection.\\
\bid{preparation-chant} & GIRM 74, referring to 48 & permitted & Preparation & No particular chant selected.\\
\bid{preface-eucharistic-prayer} & GIRM 364--365; Ordinary Time rubric & permitted & Eucharistic Prayer & Choices remain coupled and unselected; EP IV has its invariable Preface.\\
\bid{communion-psalm-or-mark} & Week XXIX, explicit alternatives & appointed alternative & Communion & Both explained; neither locally selected.\\
\bid{communion-missal-or-permitted-chant} & GIRM 87--88 & permitted & Communion & Missal routes studied; no other chant or post-Communion song selected.\\
\bid{concluding-blessing} & GIRM 90 & permitted & Conclusion & No local solemn blessing selected.\\
\end{branchtable}

The Missal Entrance and Communion antiphons are recited according to GIRM 48 and 87 when there is no singing. The Gloria and Creed are required Sunday Ordinary texts (GIRM 53,68). No Sequence, proper Offertory antiphon, uniquely appointed Preface, special Eucharistic Prayer insertion or distinctive concluding formula appears in this Week XXIX formulary. An unresolved open selection is not an additional appointed text.

\section{Appendix: Scope and Qualifications}
This study treats all ten ritual positions and eleven element keys, including both Communion alternatives. The two whole-formulary readings are editorial theological syntheses. Each Father or saint is cited for his own passage: direct exegesis, a Synoptic parallel or the identified doctrinal treatment. The prayers and readings have distinct textual histories; the independent apostolic course and shared Missal prayers retain their own roles.

The Scripture printed here is public-domain Douay--Rheims/Challoner study English, not the approved English proclaimed in the celebration. The U.S. Lectionary controls proclamation. Complete verse envelopes are explicitly distinguished from short appointed portions. Missal antiphons and prayers retain rights-safe locators; their approved English and Latin bodies are not reproduced or reconstructed. Neither the explanations nor the study passages are substitutes for liturgical recitation.

Common Missal verification belongs to the Week 29 source owner. Its inspected witnesses include the secondary digital Latin third edition of 2002, the official-hosted ICEL Antiphonary, approved FDLC extracts and publisher corroboration. Exact 2008 Latin and publisher-specific 2011 U.S. altar-book pages remain uninspected. The bounded 2008 variation check does not establish exhaustive word-for-word identity. These limits permit the present citation-led analytical treatment, not a claim of complete altar-book collation.

The reception search covers the named direct commentaries and the bounded source families recorded in the research. Bede's Mark passage was checked against PL 92, cols.236A--B; the complete Latin transcription is not thereby critically collated. Jerome's DCO witness is modern typesetting with no established historical exemplar. Corpus Thomisticum supplies identified electronic texts rather than inspected historical printings. The NPNF texts are historical translations, with their editorial notes and Augustine abridgment distinguished from the authors' own words. No claim of exhaustive patristic or modern scholarship is made.

Biblical chronology retains the alternatives and unresolved answers of the canonical chronology record. The historical appendix distinguishes source frameworks and supplies the evidence behind the concise companion's orientation. Work-wide source and rights limits, the exact bindings, and the interpretation comparisons remain in the canonical research records. This is a study aid, not an official liturgical book or an ecclesiastical approval.

\section{References}
\begin{description}[style=nextline,leftmargin=0pt,labelsep=0pt,itemsep=.35em]
\item[Scripture and liturgical books]
Douay--Rheims Bible, Challoner revision, Project Gutenberg 1581 text as registered in the source library. Isaiah 44--45; Psalms 16,32,95 (Vulgate); Matthew 21--22; Mark 10; Philippians 2; 1 Thessalonians 1,3. Roman Missal and U.S. Lectionary editions identified above; \work{Ordo lectionum Missae}, second typical edition (1981), no.145 A, pp.79--80. \work{General Introduction to the Lectionary}, §§65--69,103--104. U.S. GIRM, §§48,51,53--54,61--63,65--71,74,87--90,364--365.
\item[Dated U.S. occurrence]
USCCB, \work{Liturgical Calendar for the Dioceses of the United States of America 2026}, pp.5,41; \href{https://bible.usccb.org/bible/readings/101826.cfm}{readings for 18 October 2026}, Lectionary 145. Official online controls checked 13 October 2026.
\item[Augustine]
\work{Expositions on the Psalms}, XVI §§1--14; XXXII §§1--22; XCV §§1--15. NPNF, first series, VIII, historical English translation revised and abridged for that volume; CCEL text. Its delivery header gives an 1886 print basis while its preface is dated 1888. English headings XVII, XXXIII, XCVI use a different Psalm numbering. Exact sections used are cited beside the argument.
\item[Jerome]
\work{Commentarii in Isaiam}, XII on 45:1--7; DCO \work{Pars 3}, pp.126--129, Latin electronic typesetting. Historical exemplar not established. Short Latin terms and original English explanation only.
\item[John Chrysostom]
\work{Homilies on Matthew}, LXV §§4--5 and LXX §§1--2, Prevost translation revised by M.B. Riddle, NPNF I.10 (1888); \work{Homilies on 1 Thessalonians}, I, and \work{Homilies on Philippians}, VIII, NPNF I.13 (1889), Oxford translation revised for that volume, CCEL texts.
\item[Bede]
\work{In Marci evangelium expositio}, III on Mark 10:42--45, PL 92 (Paris, 1862), cols.236A--B; Latin Wikisource/Migne transcription, decisive passage collated with the printed image. No modern English translation used.
\item[Thomas Aquinas]
Corpus Thomisticum Latin: \work{Super Psalmos} 16 §§3--4 and 32 §§15--17; \work{Expositio super Isaiam ad litteram} 45 (85005); \work{Super Matthaeum} 22.2 (87389); \work{Super I Thessalonicenses} 1.1 (87862); \work{Super Philippenses} 2.3--4 (87838--87839). \work{Summa theologiae} III, q.48, a.4 and a.5 opening, English Dominican translation, Gutenberg 19950. The electronic heading's ``Art.3'' beneath ``Fourth Article'' is a transcription error; the structural fourth article controls the citation.
\item[Gregory Nazianzen and John Paul II]
Gregory, \work{Oration} XLV §§21--24, Browne--Swallow translation, NPNF II.7 (1894), CCEL text. John Paul II, \href{https://www.vatican.va/content/john-paul-ii/en/audiences/2002/documents/hf_jp-ii_aud_20020918.html}{General Audience, 18 September 2002}, §§1--5; original paraphrase of the modern official English.
\item[Historical orientation]
The \work{Catholic Encyclopedia} articles by Corbett, Gerard, Howlett, Souvay, Meistermann, Jacquier, Durand, MacRory, Vander Heeren, Prat and Drum are identified at their claims in the historical appendix. PBC, responsa of 26 June 1912, \work{AAS} 4 (1912), pp.463--465. USCCB, NABRE introductions to Psalms and Isaiah, official web deliveries read in the research of 13 October 2026. Historical source judgments and modern critical horizons remain distinct.
\end{description}
